% =====================================================================
% The Next Generation of Software Testing: A Critical Review of
% Agentic AI, Self-Adaptive Test Frameworks, and Autonomous Quality
% Engineering
% Template: Elsevier CAS single-column ("cas-sc"), as required by
% Computer Science Review (els-cas-templates, CTAN). Verified against
% the official els-cas-templates distribution (cas-sc.cls, identical
% by hash) and the journal's Guide for Authors: 12pt review layout,
% numbered citation style (natbib[numbers] + elsarticle-num.bst),
% Highlights (<= 85 chars/bullet), Keywords, CRediT authorship, and
% Declarations (Funding, Competing interests, Generative AI,
% Data availability) sections below.
% =====================================================================
\documentclass[a4paper,fleqn,review,12pt]{cas-sc}

\usepackage[numbers,sort&compress]{natbib}
\usepackage{graphicx}
\usepackage{amsmath,amssymb}
\usepackage{array}
\usepackage{booktabs}
\usepackage{url}
\usepackage{hyperref}
\usepackage{textcomp}
\usepackage{xcolor}
\usepackage{pgfplots}
\pgfplotsset{compat=1.18}

\hyphenation{soft-ware test-ing}

\begin{document}

\let\WriteBookmarks\relax
\def\floatpagepagefraction{1}
\def\textpagefraction{.001}

\shorttitle{The Next Generation of Software Testing}
\shortauthors{Author}

\title[mode=title]{The Next Generation of Software Testing: A Critical Review of
  Agentic AI, Self-Adaptive Test Frameworks, and Autonomous Quality Engineering}

\author{Author}
\cormark[1]
\ead{author@example.com}
\credit{Conceptualization, Methodology, Investigation, Data Curation, Formal Analysis, Writing -- original draft, Writing -- review \& editing}
\affiliation{organization={Independent Researcher},
            addressline={}, city={}, postcode={}, state={}, country={}}

\cortext[1]{Corresponding author}

\begin{abstract}
Large language models (LLMs) are increasingly embedded in autonomous and
self-adaptive software-testing pipelines, spanning unit-test generation,
graphical user interface (GUI) exploration, bug reproduction, flaky-test triage,
fuzzing, and representational state transfer (REST) API testing, promising to
automate activities long considered to require human judgment. Despite rapid
growth, the literature is fragmented across techniques, benchmarks, and claims
of effectiveness, with several influential results later shown to be overstated
once rigorously re-evaluated, making it difficult to assess the genuine state
of agentic AI in quality engineering. This systematic literature review (SLR),
following Kitchenham and PRISMA guidelines, screens 149 candidate studies drawn
from topic-scoped and protocol database searches (IEEE Xplore, ACM Digital
Library, Scopus, and Web of Science) spanning \textit{2018}--\textit{2026}, with
31 priority studies analyzed in depth across categories including surveys,
unit-test generation, multi-agent/agentic test generation, agentic
software-engineering foundations, bug reproduction/validation, mobile and GUI
testing agents, test oracles, continuous integration (CI) and flaky-test triage,
API testing, LLM-based fuzzing, and agent security. We find that multi-agent
architectures and coverage-guided feedback loops consistently improve test
quality over single-shot prompting, but the test-oracle problem,
cross-language/cross-scale generalization, fault-detection validity, context
management for large codebases, and agent security remain substantially
unresolved; several widely cited ``state-of-the-art'' results do not hold under
corrected, leakage-free evaluation. Advancing truly autonomous testing
therefore requires contamination-resistant benchmarks, fault-detection-centric
(not coverage-centric) evaluation, and dedicated attention to agent robustness,
security, and developer trust, alongside continued architectural innovation;
we outline a corresponding research agenda.
\end{abstract}

\begin{highlights}
\item Systematic review of agentic AI in software testing across 149 studies.
\item 18-category taxonomy from unit test generation to self-healing agents.
\item Synthesis of QE safety: oracle circularity, self-healing, and security.
\item Four candidate safety mechanisms proposed for autonomous test agents.
\item Ten-item research roadmap for reaching high-autonomy quality engineering.
\end{highlights}

\begin{keywords}
software testing \sep test automation \sep agentic AI \sep large language models \sep
self-healing test automation \sep systematic literature review \sep software quality assurance
\end{keywords}

\maketitle

\section{Introduction}
\label{sec:introduction}
% Introduction
%
% Suggested flow for a CS review/survey paper:
% 1. Broad context: why software testing matters
% 2. Narrowing: what is changing ("next generation")
% 3. Problem statement / gap this review addresses
% 4. Research questions (RQs)
% 5. Contributions (bulleted)
% 6. Paper structure (roadmap)

\IEEEPARstart{S}{oftware} testing remains one of the most resource-intensive
activities in the software development lifecycle, commonly consuming
30--50\% of overall project effort~\cite{placeholder2024}. Traditional
testing techniques -- manual test case design, coverage-driven unit
testing, and scripted regression suites -- are increasingly challenged
by the scale, complexity, and release velocity of modern software
systems, including microservices, machine-learning-enabled components,
and continuously deployed cloud-native applications.

In response, a new generation of testing approaches has emerged,
leveraging artificial intelligence (AI), machine learning (ML),
large language models (LLMs), self-adaptive test oracles, and
autonomous exploratory agents. These approaches promise to reduce
manual effort, improve defect detection, and adapt testing strategies
dynamically. However, the research landscape is fragmented: studies
differ in terminology (e.g., ``intelligent testing,'' ``autonomous
testing,'' ``AI-augmented testing''), evaluation benchmarks, and
maturity levels, making it difficult for researchers and practitioners
to assess the state of the art.

\subsection{Problem Statement}
No comprehensive, up-to-date systematic review consolidates these
next-generation software testing techniques under a unified
classification framework with rigorous methodological screening.

\subsection{Research Questions}
This review addresses the following research questions (RQs):
\begin{itemize}
  \item \textbf{RQ1:} What next-generation software testing techniques
        have been proposed in the literature, and how can they be
        classified?
  \item \textbf{RQ2:} What levels of automation and autonomy do these
        techniques achieve across the test life cycle (design,
        generation, execution, oracle, maintenance)?
  \item \textbf{RQ3:} What empirical evidence exists regarding the
        effectiveness, efficiency, and industrial adoption of these
        techniques?
  \item \textbf{RQ4:} What are the open challenges and research gaps
        for next-generation software testing?
\end{itemize}

\subsection{Contributions}
\begin{itemize}
  \item A systematic, reproducible literature review protocol
        (Section~\ref{sec:methodology}) covering \textit{[DATE\_RANGE]}.
  \item A taxonomy of next-generation software testing techniques
        derived from \textit{[N\_FINAL]} primary studies.
  \item A synthesis of empirical evidence on effectiveness and
        adoption barriers.
  \item A forward-looking research agenda (Section~\ref{sec:conclusion}).
\end{itemize}

\subsection{Paper Structure}
The remainder of this paper is organized as follows.
Section~\ref{sec:background} provides background and motivation.
Section~\ref{sec:methodology} details the review protocol.
Section~\ref{sec:related-work} discusses related reviews and surveys.
Section~\ref{sec:results} presents the synthesis of findings.
Section~\ref{sec:discussion} discusses implications.
Section~\ref{sec:validity} addresses threats to validity, and
Section~\ref{sec:conclusion} concludes with future directions.


\section{Background and Motivation}
\label{sec:background}
Understanding why ``agentic AI for software testing'' has become a
distinct and fast-moving research area requires situating it against
the longer trajectory of test-automation research, clarifying the
specific technical properties of large language models (LLMs) and
agentic architectures that differentiate this generation of tools from
earlier automation, and fixing a shared vocabulary for terms --
autonomy level, self-healing, the test-oracle problem, flakiness --
that recur throughout the taxonomy and synthesis in later sections.
This section does not attempt a comparative analysis of prior surveys
(that is the task of Section~\ref{sec:related-work}); instead, it
supplies the conceptual scaffolding a reader needs before engaging with
the review's results.

\subsection{A Short History of Software Test Automation}
Software testing automation has evolved through several overlapping
eras, each lowering the manual effort required to exercise a system
while raising new questions about what ``coverage'' or ``correctness''
actually guarantees. The earliest automation targeted the mechanics of
test \emph{execution} rather than test \emph{design}: capture/replay
tools recorded a tester's interactions with a GUI or API and replayed
them as scripted regression suites. This approach is cheap to adopt but
notoriously brittle -- a single UI change invalidates every recorded
locator, producing the maintenance burden that decades later motivates
self-healing locator research~\cite{h2nd,rohamo2024robotframework}.
A second era shifted the unit of automation from whole user journeys to
isolated code units, popularized by xUnit-style frameworks that made
writing and running assertions over individual functions or classes a
routine developer activity rather than a separate QA discipline; this
coverage-driven, developer-authored style of unit testing remains the
substrate that most LLM-based test-generation work still targets and
evaluates against~\cite{slr0072025,slr0082025}.

A third era pursued automated test \emph{generation} rather than mere
execution, using the structure of the program itself as a generation
signal. Search-based software testing framed test-input selection as
an optimization problem over a fitness landscape (typically branch or
line coverage), while symbolic and concolic execution reasoned about
path constraints to derive inputs that reach specific program
locations; fuzzing complemented both by mutating or generating inputs
to find crashes and robustness violations without requiring a precise
specification of correctness. These coverage-driven and property-driven
generators were effective at producing inputs but persistently weak at
producing \emph{meaningful} test oracles and at escaping coverage
plateaus once a program's branching structure became deep or
semantically opaque -- a limitation that later motivated hybridizing
search-based generation with LLMs, as in CodaMosa's strategy of
invoking a pre-trained language model to escape stalled coverage
exploration~\cite{slr0122023}. Classical fuzzers and symbolic engines
also had no access to natural-language context (requirements,
identifiers, comments, commit messages) and so could not connect a
generated input to a domain-level notion of intended
behavior~\cite{xia2024,deng2023,deng2023b}.

The most recent era layers machine learning, and subsequently large
language models and agentic orchestration, on top of this stack.
Early ML-for-testing work used learned models for narrower sub-problems
-- prioritizing test cases, predicting flaky outcomes, or ranking
likely-buggy locations -- without altering the fundamentally
human-authored nature of test design. The arrival of LLMs trained on
large code and natural-language corpora changed the unit of automation
again: a model can now read a function, its surrounding repository
context, and an informal description of intent, and directly emit a
syntactically valid, semantically plausible unit test, mutation, fuzz
target, or GUI action sequence~\cite{slr0012024,slr0022024,slr0032023}.
Agentic AI represents a further step beyond single-shot generation: a
testing agent couples an LLM with tool use (running a compiler,
executing a test, reading a stack trace, querying a repository) inside
an iterative plan-act-observe loop, allowing it to generate a test,
execute it, diagnose a failure, and revise its own output without
further human intervention~\cite{slr0232024,slr0252024,slr0242024}.
Multi-agent variants extend this further still, assigning specialized
roles (generator, critic, oracle-repairer) to separate LLM instances
that debate or vote on a shared artifact before it is
accepted~\cite{slr0192025,xu2025candor,slr0222026}. This progression --
from replaying what a human already did, to generating code that
satisfies a structural criterion, to agents that plan, act, and revise
toward a goal -- is the backdrop against which the remainder of this
review classifies next-generation testing techniques.

\subsection{Why LLMs and Agentic AI Changed the Picture}
Four technical properties of LLM-based and agentic systems explain why
this generation of tools is treated as a qualitative departure from
earlier automation rather than an incremental improvement to it.

\textbf{In-context code understanding.} Unlike symbolic or
coverage-driven generators, which reason purely over a program's
syntax or control-flow graph, LLMs are trained on vast corpora of code
paired with natural-language artifacts (comments, documentation,
commit messages, issue discussions), giving them an implicit model of
what a function is \emph{supposed} to do, not merely what paths it
contains. This lets a model propose plausible test oracles -- expected
return values, exception conditions, invariants -- directly from a
function's signature, docstring, and surrounding usage, something
earlier generators could only approximate through heuristics such as
implicit specifications mined from existing assertions or crash
oracles~\cite{dinella2022,hossain2023,o32024}.

\textbf{Natural-language-to-test generation.} Because LLMs operate
over natural language and code in a shared representation, a
requirement, bug report, or user story can itself become the
generation prompt, collapsing the traditional pipeline in which a
human first translates a requirement into a formal test specification.
Tools such as RESTGPT for API testing, NaviQAte and WebTestPilot for
web navigation, and DroidAgent and GPTDroid for mobile GUI exploration
all exploit this property to drive functional test generation directly
from natural-language intent or app-state descriptions rather than
from pre-written scripts~\cite{kim2024,shahbandeh2024,g62026,yoon2023,liu2024}.

\textbf{Tool use and multi-step, multi-agent planning loops.} Agentic
frameworks equip an LLM with callable tools -- a shell, a debugger, a
test runner, a code-search index -- and an explicit control loop that
interleaves reasoning with action, observation, and re-planning. This
allows an agent to validate its own output empirically (does the
generated test compile? does it pass? does it reproduce the reported
bug?) rather than emitting a single unverified artifact, and to
decompose a task such as ``reproduce and fix this failing CI run''
across multiple reasoning steps and, increasingly, multiple
specialized agents~\cite{slr0232024,slr0282024,slr0602026,stennett2025}.
This loop is the architectural feature that most clearly separates
agentic testing tools from earlier one-shot neural generators, and it
is also the source of the review's later distinctions between
single-agent, tool-augmented, and multi-agent architectures (RQ2).

\textbf{Self-healing and self-adaptive test maintenance.} Because an
LLM can inspect a changed UI, API schema, or DOM tree and propose an
updated locator, assertion, or oracle in place of a now-broken one,
test \emph{maintenance} -- traditionally the single largest recurring
cost of scripted and record/replay automation -- becomes an operation
the system can perform on itself. Self-healing locator frameworks for
tools such as Robot Framework and Selenium, and runtime-testing work on
self-adaptive systems more broadly, exemplify this
shift~\cite{h2nd,rohamo2024robotframework,ijam2025healenium,joseph2026zerocost,weyns2019,wohlrab2021}.
Taken together, these four properties explain both the enthusiasm
surrounding agentic testing and the specific new risks it introduces,
since a system that can rewrite its own oracles and locators can also
silently rewrite away a genuine regression.

\subsection{Key Terminology and Concepts}
The results and discussion sections rely on a small set of terms that
are used inconsistently across the primary literature; this review
adopts the following working definitions.

\textbf{Agentic AI / autonomous testing agent.} We use \emph{agentic
AI} to denote an LLM-based system that perceives a software artifact
or execution environment, forms a plan, invokes external tools to act
on that environment, observes the result, and iterates -- as opposed to
a single-pass model invocation that maps an input prompt directly to
an output artifact~\cite{slr0042024,slr0052024,slr0062024}. An
\emph{autonomous testing agent} is an agentic system whose task is
specifically test design, generation, execution, triage, or
maintenance, potentially spanning several of these life-cycle stages
without human intervention between them~\cite{slr0882020,slr0662026}.

\textbf{Self-healing test automation.} This refers to mechanisms --
typically, but not exclusively, LLM-based -- that automatically detect
a broken element of a test (a UI locator, an API contract assumption,
an assertion that no longer matches intended behavior) and repair it at
or near runtime, without requiring a human to re-author the affected
test~\cite{h2nd,rohamo2024robotframework,ijam2025healenium}. Self-healing
is a special case of the broader notion of a self-adaptive system,
which continuously monitors and adjusts its own structure or behavior
in response to a changing environment~\cite{weyns2019,hnsel2018,sa42021}.

\textbf{The test-oracle problem.} The test-oracle problem is the
long-standing difficulty of deciding, for a given test input, what the
\emph{correct} output or behavior should be, absent a complete and
trustworthy specification. Neural and LLM-based oracle generation
(e.g., TOGA, TOGLL, AugmenTest) attempts to approximate this judgment
by learning from existing assertions, exception patterns, or developer
intent expressed in code and
comments~\cite{dinella2022,o32024,o42025,bodicoat2026oracle,molina2024oracleautomation}.
Using an LLM to both generate a test and judge whether its own
generated behavior is correct raises a circularity concern -- an agent
testing code it (or a closely related model) produced -- that is
compounded by documented self-preference biases in LLM-as-judge
settings and by evidence that models struggle to self-correct their own
reasoning or code without external
feedback~\cite{wataoka2024selfpreference,panickssery2024evaluators,huang2023selfcorrect,olausson2023selfrepair}.

\textbf{Flaky tests.} A flaky test is one that can pass or fail across
repeated executions of the same code under test, without any change to
that code, typically due to non-determinism such as timing, concurrency,
test-order dependence, or external-resource instability. Flaky tests
undermine the trustworthiness of continuous-integration signals and
have motivated a substantial body of machine-learning and, more
recently, LLM-based classification and repair
work~\cite{fatima2025,alshammari2024flakyclassifiers,parry2023cannier,leinen2026detectedundetected,parry2025systemic,lam2020lifecycle,meta2020pfs}.
Flakiness is directly relevant to self-healing test automation because
a system that automatically repairs a ``failing'' test must first
distinguish a genuine regression from an artifact of flakiness, or risk
silently suppressing real defects.

\textbf{Autonomy-level framing (L0--L5).} To compare how much human
oversight a given technique still requires, this review adopts the
autonomy-level scale used later in the roadmap and discussion
sections, following the general levels-of-autonomy framing for AI
agents of \cite{feng2025autonomylevels}: autonomy ranges from no
agentic behavior (a human performs the task with no AI assistance)
through assistive and co-pilot levels (the system suggests artifacts a
human reviews and approves) to conditional and high autonomy (the
system executes multi-step plans with narrow human checkpoints) and,
at the upper end, full autonomy, in which a testing agent designs,
executes, triages, and maintains tests across a life cycle with no
routine human intervention. This scale, together with the ISO/IEC
guidance on testing AI-based systems and on testing AI more
generally~\cite{iso2020tr29119,iso2025ts42119}, gives the taxonomy in
Section~\ref{sec:results} a consistent axis along which to classify
otherwise heterogeneous techniques (RQ2).

\subsection{Motivation: Why This Topic Needs a Rigorous Synthesis Now}
Four converging pressures motivate a systematic, rather than
narrative, review of this literature at this particular moment.

\textbf{Fragmentation of claims and terminology.} The corpus assembled
for this review spans 149 candidate studies published between 2018 and
2026, distributed across at least eighteen partially overlapping
taxonomy categories -- unit-test generation, GUI and mobile testing
agents, API testing, fuzzing, mutation testing, oracle generation,
flaky-test classification, self-adaptive runtime testing, CI-failure
triage, security of coding and testing agents, and more. Authors
routinely use different labels (``intelligent testing,'' ``autonomous
testing,'' ``AI-augmented testing,'' ``agentic software engineering'')
for closely related ideas, and several existing surveys already
attempt partial consolidations of this
space~\cite{slr0012024,slr0022024,slr0032023,slr0092024,slr0102024},
underscoring that the field has grown too large and too fast for any
single narrative account to track consistently; a reproducible,
protocol-driven synthesis is needed to make comparisons across these
overlapping categories meaningful (RQ1).

\textbf{Inconsistent benchmarks and limited generalizability.}
Empirical claims in this area are frequently anchored to small,
language-specific, or project-specific benchmarks, and several
targeted replications show that headline results are sensitive to
exactly how a benchmark is constructed and evaluated. Independent
re-evaluations of flaky-test classifiers, for instance, have found that
reported accuracy drops substantially once unrealistic class balance or
weak validation protocols are
corrected~\cite{alshammari2024flakyclassifiers,parry2023cannier}, and
studies of neural test-oracle generation report similarly large gaps
between original and large-scale replicated
evaluations~\cite{hossain2023}. Without a systematic review that
extracts and compares evaluation methodology across studies -- not
merely headline numbers -- the field risks treating isolated,
possibly fragile results as generalizable evidence (RQ3).

\textbf{Data contamination and benchmark leakage.} Because LLMs are
trained on enormous, largely undisclosed corpora that plausibly include
public GitHub repositories, there is a concrete risk that a model
``generating'' a correct test for a benchmark problem is in fact
recalling memorized training data rather than reasoning about the code
under test. Direct studies of this risk -- measuring whether LLMs
memorize widely used bug benchmarks and quantifying inter-dataset code
duplication and leakage across popular training and evaluation
corpora -- show the problem is non-hypothetical, and have motivated
the construction of contamination-resistant benchmarks built from
recently created bug
datasets~\cite{slr0302025,slr0322024,slr0312023,slr0292025}. A
synthesis that explicitly tracks which primary studies address, ignore,
or are vulnerable to contamination is necessary before the field's
quantitative claims about LLM testing effectiveness can be taken at
face value.

\textbf{The hype-versus-evidence gap and the quality-engineering safety
angle.} Enthusiastic framings of agentic AI as a wholesale replacement
for human test engineering coexist with sober industrial reports that
document narrower, assistive deployments and persistent practitioner
skepticism~\cite{slr0372025,slr0382024,slr0392025,slr0832025}, and with
emerging survey and position work explicitly interrogating how far
agentic autonomy in software engineering has actually
progressed~\cite{slr0662026,slr0692026,slr0722026,slr0842025,slr0852025}.
Closing this gap matters beyond academic tidiness because autonomous
testing agents introduce safety concerns that earlier, purely
assistive tools did not: when an agent both writes production code and
generates the tests meant to validate it, the two artifacts can share
the same blind spots, producing a circularity in which a test suite
systematically fails to catch the very class of defect the generating
model is prone to introduce~\cite{xu2025candor,slr0192025}. Testing
agents are also newly exposed, tool-using software themselves, and are
therefore subject to indirect prompt injection, unsafe tool
invocation, and sandbox-escape attacks that have been demonstrated
against LLM coding and tool-augmented agents more
broadly~\cite{zhan2024injecagent,debenedetti2024agentdojo,guo2024redcode,she2025queryipi,sakib2026trustverify,kozak2025developeraid,shi2025progent,wu2024isolategpt}.
Finally, as autonomy increases toward the higher end of the L0--L5
scale, questions of accountability, overreliance, and governance --
who is responsible when a self-healing agent masks a regression, and
what standards should govern a Level-5 testing agent -- move from
speculative to
urgent~\cite{treude2026accountable,cheng2026threepillar,passi2025overreliance,ibrahim2025overreliance,kiran2026humanai}.
Together, fragmentation, benchmark fragility, contamination risk, and
this emerging safety angle constitute the gap this systematic review
is designed to address, following an SLR protocol in the tradition of
\cite{kitchenham2007} and reported per the PRISMA
guidelines~\cite{page2021prisma}.


\section{Research Methodology}
\label{sec:methodology}
% Research Methodology
%
% Follow Kitchenham & Charters (2007) SLR guidelines and/or PRISMA 2020.
% This section should be fully reproducible -- cross-reference
% literature-review/ protocol documents in the repository.

\subsection{Review Protocol}
This review follows the systematic literature review (SLR) guidelines
of Kitchenham and Charters~\cite{kitchenham2007} and reports study
selection using the PRISMA 2020 flow diagram~\cite{page2021prisma}.
The complete, machine-readable protocol (search strings, databases,
inclusion/exclusion criteria, and quality assessment checklist) is
versioned in \texttt{literature-review/} of the accompanying
repository for full reproducibility.

\subsection{Search Strategy}
\textit{[Describe digital libraries searched (e.g., IEEE Xplore, ACM
Digital Library, Scopus, Web of Science, Google Scholar), the search
string, and the time window. See
\texttt{literature-review/search-strategy.md}.]}

\subsection{Inclusion and Exclusion Criteria}
\textit{[Summarize criteria; full list in
\texttt{literature-review/inclusion-exclusion-criteria.md}.]}

\subsection{Study Selection Process}
\textit{[Describe the PRISMA stages: identification, screening,
eligibility, inclusion. Reference the PRISMA diagram generated from
\texttt{literature-review/prisma-flow.md}.]}

\subsection{Data Extraction and Synthesis}
\textit{[Describe the data extraction form fields and the synthesis
method -- e.g., narrative synthesis, thematic analysis, or
meta-analysis where applicable.]}

\subsection{Quality Assessment}
\textit{[Describe the quality assessment checklist used to score each
primary study, e.g., based on the DARE or CASP criteria adapted for
software engineering.]}


\section{Related Work}
\label{sec:related-work}
% Related Work
%
% Position this review against prior surveys/reviews in the area.
% Use a comparison table (see paper/tables/) highlighting scope,
% time period, and gaps addressed by this review.
\textit{[Summarize prior surveys and systematic reviews on software
testing, test automation, and AI for software engineering. Clearly
state how this review differs in scope, recency, or methodological
rigor. Consider including a comparison table such as
Table~\ref{tab:related-reviews}.]}

% Example table include:
% \input{tables/related_reviews_comparison.tex}


\section{Results and Analysis}
\label{sec:results}
% Results and Analysis
%
% Organize by Research Question (RQ1-RQ4). Include demographic
% analysis (publication trends by year/venue), taxonomy figures,
% and synthesis tables generated from data/processed/.

\subsection{Overview of Selected Studies}
\textit{[Report N\_FINAL primary studies, publication year
distribution, venue distribution. Reference
Figure~\ref{fig:pub-trend} generated from
\texttt{scripts/analyze\_literature.py}.]}

% \begin{figure}[t]
%   \centering
%   \includegraphics[width=\columnwidth]{figures/publication_trend.pdf}
%   \caption{Publication trend of next-generation software testing
%   studies by year.}
%   \label{fig:pub-trend}
% \end{figure}

\subsection{RQ1: Taxonomy of Techniques}
\textit{[Present the taxonomy derived from thematic coding. Consider
Table~\ref{tab:taxonomy}.]}

\subsection{RQ2: Automation and Autonomy Levels}
\textit{[Classify studies by automation level across the test
life cycle: design, generation, execution, oracle, maintenance.]}

\subsection{RQ3: Empirical Evidence}
\textit{[Synthesize reported effectiveness metrics (e.g., fault
detection rate, coverage, precision/recall of generated oracles) and
industrial case studies.]}

\subsection{RQ4: Open Challenges}
\textit{[List and discuss identified gaps: explainability,
reproducibility, benchmark standardization, flaky test handling,
cost of adoption, etc.]}


\section{Discussion}
\label{sec:discussion}
% Discussion
%
% Interpret findings, connect back to RQs, discuss implications for
% researchers, practitioners, and tool vendors.
This section interprets the results reported in
Section~\ref{sec:results} against the four research questions, then
draws out implications for three audiences -- researchers, tool
vendors, and practitioners adopting these techniques in industrial
pipelines -- before turning to the quality engineering safety synthesis
that spans all four RQs.

\subsection{Revisiting RQ1: What Techniques Exist, and How Mature Are They?}
The 18-category taxonomy shows that ``next-generation software
testing'' is not a single technique but a loosely federated family of
architectures unified mainly by their reliance on an LLM as a
reasoning component somewhere in the pipeline. At one extreme sit
single-shot code-completion-style test generators that differ from
classical automated test generation only in using a pretrained
language model instead of a symbolic or search-based engine; at the
other extreme sit genuinely agentic systems that plan a multi-step
investigation, call external tools (compilers, debuggers, browsers,
issue trackers), observe the results, and revise their own plan with
little or no human supervision. The evidence base is considerably
stronger for the former than the latter: unit-test generation and
mutation-testing techniques (categories well represented in the
corpus) benefit from a decade of prior automated-test-generation
research to build on and compare against, whereas multi-agent and
GUI/mobile testing agents are newer, smaller in study count, and more
often evaluated on bespoke benchmarks constructed by the same team
that built the tool -- a pattern that, as discussed under RQ3 in
Section~\ref{sec:results}, correlates with inflated effectiveness claims. The practical
implication for RQ1 is that technique maturity should not be inferred
from architectural sophistication: a simple, well-evaluated single-shot
generator may currently be a safer adoption choice than an elaborate
multi-agent pipeline whose claimed gains have not been independently
reproduced.

\subsection{Revisiting RQ2: What Autonomy Levels Are Actually Achieved?}
Despite frequent use of the word ``autonomous'' in paper titles, the
corpus shows that most deployed systems operate with a human firmly in
the loop at one or more critical junctures: approving generated test
code before it enters a suite, triaging flagged flaky tests rather than
having them auto-quarantined, or reviewing agent-authored pull requests
before merge. Genuinely low-supervision operation is concentrated in
narrow, well-bounded sub-tasks (e.g., locator self-healing, mutation
generation) rather than across the full test life cycle (design,
generation, execution, oracle judgment, and maintenance). This gap
between rhetorical and demonstrated autonomy is precisely why this
review adopts an explicit autonomy-level framing (L0--L5, following
\cite{feng2025autonomylevels}) in the roadmap and governance discussion
below: without a shared vocabulary for \emph{how much} supervision a
system actually requires, ``autonomous testing'' risks becoming a
marketing label rather than a measurable engineering property.

\subsection{Revisiting RQ3: How Rigorous Is the Supporting Evidence?}
The clearest cross-cutting finding under RQ3 is that benchmark design
problems are not isolated incidents but a structural feature of this
literature. Studies in the ``Benchmarks and data contamination''
category document cases where headline accuracy numbers for neural
test-oracle generation and flaky-test classification fell substantially
once training/test leakage and class-balance artifacts were corrected,
and several comparative studies found that simple, non-agentic
baselines matched or exceeded elaborate multi-agent pipelines at a
fraction of the inference cost. This matters for how the remaining
results in this review should be read: wherever a specific effectiveness
number is cited, it should be understood as a claim made \emph{under
the evaluation conditions of the original study}, not a validated,
generalizable property of the technique. For researchers, this argues
for mandatory leakage audits and cost-normalized comparisons as a
condition of publication in this area; for practitioners, it argues
for re-validating any vendor or paper's effectiveness claims against a
held-out, in-house test set before trusting them in a production
pipeline.

\subsection{Revisiting RQ4: What Remains Open?}
RQ4 is answered most directly by the quality-engineering-safety
synthesis below and by the ten-item research roadmap in
Section~\ref{sec:conclusion}, but two cross-cutting observations are
worth stating plainly here. First, the ``Human factors: trust and
acceptance'' and ``Security of coding and testing agents'' categories
are, jointly, the largest single block of the corpus (30 of 149
studies), which indicates that the research community itself already
recognizes trust and security as the dominant open problem, even though
tooling and standards have not caught up with that recognition.
Second, almost none of the reviewed studies propose or evaluate a
mechanism for \emph{detecting} when an autonomous testing agent has
gone wrong in a way that is not immediately visible as a test failure
(for example, a self-healing agent that has quietly relaxed an
assertion, or a judge model that shares a blind spot with the
generator it is judging). This asymmetry -- extensive work on making
agents capable, comparatively little on making their failures legible
-- is the throughline connecting all four quality-engineering-safety
threads discussed next.

\subsection{Implications for Practitioners and Tool Vendors}
\label{subsec:implications}
For teams evaluating these techniques for production adoption, three
practical implications follow from the synthesis presented in
Section~\ref{sec:results} and Section~\ref{sec:qe-safety}. First, treat
any claim of autonomous operation as scoped to a specific sub-task
(e.g., ``autonomous locator repair'') rather than as a property of an
entire pipeline, and verify the autonomy level against the L0--L5
framing discussed in Section~\ref{sec:conclusion} before relying on it
for unsupervised operation. Second, demand leakage-free, in-house
validation of any effectiveness claim before adoption, given the
documented frequency with which benchmark-reported numbers fail to
reproduce under corrected evaluation. Third, budget for governance
and security from the outset: indirect prompt injection against coding
and testing agents is a demonstrated, quantified threat (Section~\ref{sec:results}),
and human review catches only a fraction of agent-introduced defects
in practice, so relying on manual code review as a sole defense is
untenable.

\subsection{Quality Engineering Safety}
\label{sec:qe-safety}
A cross-cutting theme emerges across four otherwise distinct bodies of
evidence gathered for RQ4: security of testing/coding agents
\cite{zhan2024injecagent,debenedetti2024agentdojo,guo2024redcode,
she2025queryipi,sakib2026trustverify,kozak2025developeraid,
shi2025progent,wu2024isolategpt,wu2026skillshield,xie2026vibeguard,
kokkonis2025rosa}; test-oracle independence and the circularity of AI
testing AI-written code
\cite{bodicoat2026oracle,molina2024oracleautomation,xu2025candor,
wataoka2024selfpreference,panickssery2024evaluators,huang2023selfcorrect,
olausson2023selfrepair}; self-healing/flaky-test stability
\cite{alshammari2024flakyclassifiers,parry2023cannier,
leinen2026detectedundetected,parry2025systemic,lam2020lifecycle,
meta2020pfs,joseph2026zerocost,rohamo2024robotframework,ijam2025healenium};
and governance/accountability for highly autonomous quality engineering
\cite{iso2020tr29119,iso2025ts42119,feng2025autonomylevels,
treude2026accountable,cheng2026threepillar,ibm2024trew,
passi2025overreliance,ibrahim2025overreliance,kiran2026humanai}.

Reading these four threads together reveals a shared structure. In each
case, the core risk is not that an agent fails to perform a task, but
that it performs the task in a way that generates false confidence: an
oracle that quietly accepts subtle bugs because the judge model shares
pretraining blind spots with the generator; a self-healing test that
passes again by modifying a locator to point to the wrong element or by
weakening an assertion; a testing agent hijacked by indirect prompt
injection embedded in a third-party library or test fixture; or a
Level~4--5 autonomous pipeline making release-sign-off decisions without
traceable human accountability. Classical software testing research
relied on an implicit assumption of an independent, trusted, and
adversarially neutral tester; autonomous quality engineering violates
all three assumptions simultaneously.

We argue that resolving this cluster requires treating \emph{quality
engineering safety} as a distinct sub-discipline with its own
mechanisms, rather than as a set of disconnected caveats. Four
candidate mechanisms following this principle are proposed in
Section~\ref{sec:safety-framework} and the supplementary material:
a layered Trust Boundary for Test Agents (TBTA) for the security thread,
Provenance-Partitioned Oracle Verification (PPOV) for the oracle-
circularity thread, Mutation-Gated Dual-Signal Heal Verification
(MG-DSV) for the self-healing thread, and an Autonomy-Tiered Quality
Sign-off Protocol (ATQSP) for the governance thread. None of these four
mechanisms has yet been implemented or empirically evaluated in this
review; they are synthesized proposals for future work, not validated
contributions, and should be read as scoping the research agenda in
Section~\ref{sec:conclusion} rather than as established results.

A further observation, not explored in depth by any single reviewed
study but implied by reading all four threads together, is that the
``second independent signal'' principle has a recursive cost: verifying
a verifier is itself a quality-engineering problem. TBTA's sandboxing
layer must itself be monitored for escape; PPOV's provenance-partitioned
oracle must itself avoid inheriting the generator's blind spots through
shared pretraining data; MG-DSV's mutation-adequacy gate must itself be
protected from the same metric-gaming pressure that motivated the
``Benchmarks and data contamination'' findings under RQ3; and ATQSP's
sign-off protocol is only as trustworthy as the audit log it depends
on. None of the four mechanisms, as currently formalized, fully closes this
regress, which is itself evidence for treating ``verification of
verifiers'' as a standing research problem rather than a one-time
design fix -- a point made explicit in open problem~9 of the roadmap in
Section~\ref{sec:conclusion}.

\subsection{Toward a Testable Safety Framework: Safety-Assured Agentic Quality Engineering}
\label{sec:safety-framework}
The ATQSP sketch introduced in Section~\ref{sec:qe-safety} answers \emph{what} a governance artifact for
autonomous test agents should record, but it does not yet specify
\emph{how} an organization should decide, in a principled and
revisable way, how much autonomy a given agent has earned in a given
context. We elaborate ATQSP into a fuller framework, Safety-Assured
Agentic Quality Engineering (SAAQE), that borrows its central
metaphor from drug-safety practice: an agent, like a drug, is granted
a privilege (dosage, autonomy level) only as far as its accumulated,
monitored track record justifies, and that privilege is continuously
re-evaluated rather than granted once and forgotten. Framed this way,
the four mechanisms proposed in Section~\ref{sec:qe-safety} (TBTA, PPOV,
MG-DSV, and ATQSP) become the \emph{safety guardrails} that feed evidence into a
single accountability substrate, summarized in
Table~\ref{tab:saaqe-components}.

\begin{table}[t]
\caption{Components of the proposed Safety-Assured Agentic Quality
Engineering (SAAQE) framework.}
\label{tab:saaqe-components}
\begin{tabular}{p{0.3\linewidth}p{0.6\linewidth}}
\toprule
Component & What it does \\
\midrule
Safety guardrails &
Five checks: oracle independence from the code under test (PPOV),
mutation-gated self-healing safety (MG-DSV), a sandboxed security
boundary (TBTA), risk-based mandatory human review, and explicit
accountability assignment for release decisions (ATQSP). \\
Knowledge graph &
Traces each escaped defect back through the code and tests to the
specific agent(s) that wrote, healed, or approved them, using
commit-level provenance and bug-inducing-change identification. \\
Reliability score &
A Bayesian score per agent and per task context: clean work raises
the score, escaped defects lower it in proportion to severity, and a
confidence interval reflects the amount of accumulated evidence. \\
Safety boundary (``therapeutic window'') &
The autonomy range in which effort saved by the agent remains high
while the escaped-defect rate stays within an organization-defined
acceptable bound. \\
Privileging &
An agent is granted autonomy level L0--L5
\cite{feng2025autonomylevels} only once its reliability score clears
that level's threshold, and is automatically demoted a level when the
score falls below it -- operationalizing ATQSP's Autonomy Certificate
as a continuously re-evaluated, rather than one-time, grant. \\
Phased rollout and monitoring &
Sandbox evaluation, then shadow mode (the agent acts but a human
retains the decision), then limited production rollout, then full
production, with continuous monitoring of the escaped-defect rate at
each stage before advancing. \\
\bottomrule
\end{tabular}
\end{table}

Read together with the roadmap items in Section~\ref{sec:conclusion}, SAAQE's intended
contribution is fourfold: (i) the framework itself, unifying the four
guardrail mechanisms under one accountability substrate; (ii) a small
set of safety metrics -- escaped-defect rate, masked-regression rate
(valid test behavior silently weakened by a heal or an oracle change),
and attack-success rate under adversarial prompt injection -- that
make ``safety'' measurable rather than aspirational; (iii) a
benchmark suite assembled from existing public datasets (proposed
next); and (iv) an open-source reference implementation of the
knowledge graph and reliability-scoring components, so that the
scoring mechanism itself is auditable rather than a proprietary black
box. We position this contribution within the niche of
\emph{safety-assured agentic quality engineering}: unlike the
capability-focused benchmarks surveyed under RQ1--RQ3, SAAQE is
explicitly a safety-governance layer that sits above whichever
test-generation or self-healing technique an organization already
uses, and is therefore complementary to, not competitive with, the
techniques cataloged in the taxonomy (Section~\ref{sec:results}).

\subsection{An Empirical Validation Plan Using Public Datasets}
\label{sec:saaqe-validation}
Like the four guardrail mechanisms it builds on, SAAQE is a proposed,
not yet empirically validated, mechanism; this subsection outlines how
it \emph{could} be validated before any claim of effectiveness is
made, using data that -- unlike an individual organization's private
CI/CD and defect-tracker history -- is public and therefore
independently reproducible. The guiding idea is to \emph{replay
history}: score agents retrospectively on past work using only
information available at the time, and test whether the resulting
score predicts defects that escaped and were discovered later.
Table~\ref{tab:saaqe-data} summarizes candidate public datasets and
the specific SAAQE component each would exercise.

\begin{table}[t]
\caption{Candidate public datasets for validating the proposed SAAQE
framework.}
\label{tab:saaqe-data}
\begin{tabular}{p{0.26\linewidth}p{0.3\linewidth}p{0.34\linewidth}}
\toprule
Data & What it gives & Use it to test \\
\midrule
AIDev \cite{li2026aidev} &
932{,}791 real agent-authored pull requests from five coding agents
across 116{,}211 repositories &
Reliability scores and the knowledge graph on real, in-the-wild agent
work \\
GitHub issues and bug-fix commits traced with SZZ
\cite{sliwerski2005szz} &
Which earlier commits introduced each later-reported bug &
Escaped defects (bugs found after merge), linked back to the
responsible agent \\
SWE-bench \cite{jimenez2024swebench}, SWT-Bench \cite{slr0282024},
TestGenEval \cite{jain2024testgeneval} &
Agent results on real GitHub issues and tests &
Score calibration; oracle-independence guardrail (PPOV) \\
Defects4J \cite{just2014defects4j}, BugsInPy
\cite{widyasari2020bugsinpy} &
Known, real, previously-fixed faults (prefer the newer, less
contamination-prone releases) &
Whether agent-written tests catch known faults; the masked-regression
rate \\
Flaky-test datasets (e.g., as used by FlakyLens
\cite{slr0442025} and FlakyGuard \cite{slr0492025}) &
Real flaky tests and their fixes &
Heal-safety guardrail (MG-DSV) \\
AgentDojo \cite{debenedetti2024agentdojo}, RepoGuardBench
\cite{li2026repoguardbench} &
Prompt-injection attacks against coding/testing agents &
Security-sandbox guardrail (TBTA) \\
Organization-internal CI/CD and defect-tracker history (where access
permits) &
Production escaped defects with real severity labels &
The strongest validation, as an industrial case study \\
\bottomrule
\end{tabular}
\end{table}

A four-step validation protocol follows from this data: (1)
\emph{build the knowledge graph} from AIDev by linking each agent's
pull requests to the files, tests, and later bug-fix commits
identified via SZZ; (2) \emph{replay the reliability score in time
order}, updating each agent's score at each historical time step using
only evidence that would have been available at that time, to avoid
leaking future information into a retrospective score; (3)
\emph{check predictive validity}, measuring with AUC, precision, and
recall whether agents whose replayed score was low at time $t$
produced more escaped defects in the subsequent period than agents
whose score was high; and (4) \emph{test the safety boundary
counterfactually}, estimating how many of the actually-escaped defects
in the historical record would have been prevented, and how much
measured productivity would have been forgone, had autonomy been
restricted to agents whose score cleared the proposed threshold at
the time. A successful outcome on all four steps would support three
claims: \emph{validity} (the reliability score predicts future
escaped defects), \emph{value} (the safety boundary reduces escaped
defects at an acceptable productivity cost), and \emph{generalization}
(the relationship holds across repositories, languages, and agents).
Three limitations should be stated up front, consistent with this
review's broader emphasis on methodological caution (RQ3,
Section~\ref{sec:threats}): agent-authored pull requests are a recent
phenomenon, so the available history for replay is short;
SZZ-identified bug-inducing commits are known to include labeling
errors, so any SZZ-derived escaped-defect count should be reported
alongside an estimate of SZZ's own error rate rather than treated as
ground truth; and public repositories rarely carry official severity
ratings, so severity must be estimated from proxies such as issue
labels and reporter language, a construct-validity limitation that
should be disclosed rather than hidden behind a single aggregate
metric.


\section{Threats to Validity}
\label{sec:threats}
\label{sec:validity}
% Threats to Validity
%
% Standard SLR validity dimensions: construct, internal, external,
% conclusion (following Wohlin et al. or Runeson & Host).
\subsection{Construct Validity}
\textit{[E.g., search string coverage, terminology ambiguity.]}

\subsection{Internal Validity}
\textit{[E.g., researcher bias in screening/coding, mitigated via
dual-reviewer agreement -- report Cohen's kappa if applicable.]}

\subsection{External Validity}
\textit{[E.g., generalizability of findings beyond the sampled
venues/databases/time window.]}

\subsection{Conclusion Validity}
\textit{[E.g., reliability of the data extraction and synthesis
process.]}


\section{Conclusion and Future Work}
\label{sec:conclusion}
% Conclusion and Future Work
%
% 1. Summary of review findings (answering RQ1-RQ4 briefly)
% 2. Synthesis of implications
% 3. Research roadmap: prioritized open problems
% 4. Concluding remarks

\subsection{Summary of Findings}
This systematic literature review has synthesized 149 primary
studies (2018--2026) on next-generation software testing, spanning
large-language-model-based test generation, multi-agent architectures,
self-adaptive and self-healing test automation, test oracles, and
autonomous quality engineering. We answered four research questions:

\begin{itemize}
  \item \textbf{RQ1 (Taxonomy):} We identified 18 technique
        categories across the test life cycle, from unit test
        generation to multi-agent systems and self-healing automation.
  \item \textbf{RQ2 (Autonomy):} Most techniques operate at
        human-in-the-loop assistance levels (L1--L3); true agentic
        autonomy (L4--L5) is concentrated in narrow sub-tasks such as
        GUI exploration and flaky-test triage, with significant human
        supervision remaining for test design and oracle validation.
  \item \textbf{RQ3 (Evidence):} Empirical evidence demonstrates
        consistent coverage gains from LLM-based and multi-agent
        generation, but fault-detection effectiveness is less rigorously
        measured, and several widely cited findings do not hold under
        corrected, leakage-free evaluation.
  \item \textbf{RQ4 (Challenges):} Key open problems include the
        test-oracle problem, benchmark contamination,
        cross-language/cross-scale generalization, context management
        for large codebases, and the security and trustworthiness of
        autonomous testing agents.
\end{itemize}

Building on the primary synthesis, this review extended its scope to
address four emergent \emph{quality engineering safety} threads:
security of testing/coding agents, test-oracle independence and the
circularity of AI testing AI-written code, self-healing/flaky-test
stability, and governance/accountability for high levels of autonomy.
This extension was driven by an empirical finding in the corpus itself:
nearly a third of the categorized primary studies address risk and
trust rather than capability. The resulting synthesis formalizes
several of these RQ4 gaps as a connected set of ten open problems and
four candidate, as-yet-unvalidated safety mechanisms.

\subsection{Research Roadmap: Ten Open Problems}
\label{sec:research-roadmap}
Building on the gaps surfaced throughout the synthesis
(Section~\ref{sec:discussion}), we propose the following research
roadmap for next-generation, autonomous software testing. Each item is
stated as an open problem rather than a solution, to orient future
empirical and methodological work (RQ4). Items 1, 4, 7, 8, 9, and 10 —
collectively our \emph{Quality Engineering Safety} thread — are backed
by a dedicated literature synthesis and a candidate safety mechanism each
(detailed in Section~\ref{sec:discussion}; corresponding primary citations
are included in the bibliography):

\begin{enumerate}
  \item \textbf{Oracles independent of the code under test.} Most
        LLM-based oracles are trained or prompted on artifacts derived
        from the same codebase they judge, risking circular validation;
        open problem: oracle construction strategies that do not
        depend on the implementation being tested. Empirical evidence
        of shared generator/judge failure modes is reported by
        \cite{panickssery2024evaluators,wataoka2024selfpreference,
        huang2023selfcorrect,olausson2023selfrepair}; partial mitigations
        (consensus panels, context grounding) are evaluated by
        \cite{xu2025candor,bodicoat2026oracle,molina2024oracleautomation}.
        We propose \emph{Provenance-Partitioned Oracle Verification
        (PPOV)} as a candidate mechanism (formalized in
        Section~\ref{sec:safety-framework}).
  \item \textbf{Probabilistic guarantees on test-suite adequacy.}
        Coverage metrics do not bound fault-detection probability for
        LLM-generated suites; open problem: statistically grounded
        adequacy criteria with confidence guarantees.
  \item \textbf{Contamination-free fault-detection benchmarks.}
        Existing benchmarks risk training-data leakage, inflating
        reported accuracy; open problem: benchmarks that measure true
        fault detection and are verifiably free of contamination.
  \item \textbf{Stability of self-healing loops.} Self-healing locators
        and tests can silently suppress genuine regressions while
        recovering pass rates; open problem: bounding and certifying the
        correctness of automated repairs under distributional UI and
        API drift. Empirical evidence on flaky-classifier limits is
        reported by \cite{alshammari2024flakyclassifiers,parry2023cannier,
        leinen2026detectedundetected}, systemic flakiness by
        \cite{parry2025systemic}, fix stability by \cite{lam2020lifecycle},
        and continuous scoring by \cite{meta2020pfs}. We propose
        \emph{Mutation-Gated Dual-Signal Heal Verification (MG-DSV)} as
        a candidate mechanism (formalized in
        Section~\ref{sec:safety-framework}).
  \item \textbf{Cross-language and cross-ecosystem generalization.}
        Research is heavily concentrated in Java and Python; open
        problem: techniques and benchmarks spanning systems languages,
        functional paradigms, and embedded platforms.
  \item \textbf{Scalable context management for repository-level testing.}
        Single-file/method context limits LLM reasoning; open problem:
        hierarchical retrieval, code-graph representations, and memory
        architectures for whole-repository test generation and triage.
  \item \textbf{Security and integrity of testing agents.} Testing agents
        execute code with elevated privileges, creating attack vectors
        (prompt injection, code poisoning, exfiltration); open problem:
        least-privilege execution environments, verified tool calling,
        and supply-chain defenses tailored to testing pipelines. Attack
        evidence is provided by
        \cite{zhan2024injecagent,debenedetti2024agentdojo,guo2024redcode,
        she2025queryipi,sakib2026trustverify,kozak2025developeraid}, with
        candidate defenses by
        \cite{shi2025progent,wu2024isolategpt,wu2026skillshield,
        xie2026vibeguard,kokkonis2025rosa}. We propose a layered
        \emph{Trust Boundary for Test Agents (TBTA)} architecture
        (formalized in Section~\ref{sec:safety-framework}).
  \item \textbf{Human--agent collaboration.} When to escalate to a
        human, who is accountable for agent-authored tests, and how
        trust is calibrated remain open; open problem: escalation and
        accountability frameworks for mixed human/agent QE teams.
        Evidence on overreliance and trust calibration is reported by
        \cite{passi2025overreliance,ibrahim2025overreliance,ibm2024trew},
        with QA-specific escalation proposals in \cite{kiran2026humanai}.
  \item \textbf{The circularity problem of AI testing AI-written code.}
        As AI increasingly writes both code and its tests, independence
        between implementation and verification erodes; open problem:
        validation strategies that preserve independence when both
        artifacts share a generative origin. See the oracle-
        independence evidence under item 1 in this roadmap
        \cite{panickssery2024evaluators,wataoka2024selfpreference},
        which is the same underlying mechanism.
  \item \textbf{Governance standards for Level-5 (fully autonomous)
        testing.} Existing standards (e.g., ISO/IEC/IEEE~29119) do not
        address fully autonomous test agents; open problem: extending
        or replacing such standards with governance criteria for L5
        autonomy. The current standards baseline is
        \cite{iso2020tr29119,iso2025ts42119}; autonomy-level and
        accountability proposals that could extend it are given by
        \cite{feng2025autonomylevels,treude2026accountable,
        cheng2026threepillar}. We propose an \emph{Autonomy-Tiered
        Quality Sign-off Protocol (ATQSP)} as a candidate extension,
        elaborated in Section~\ref{sec:safety-framework} into the fuller
        Safety-Assured Agentic Quality Engineering (SAAQE) framework,
        with a proposed public-dataset validation protocol in
        Section~\ref{sec:saaqe-validation}.
\end{enumerate}

These ten items map onto, and partially answer, recurring gaps
identified across the included studies (Section~\ref{sec:discussion}):
test-oracle reliability and independence, unmeasured fault-detection
adequacy, benchmark/data contamination, self-healing/flaky-test
stability, reproducibility of agentic results, and security/governance
of autonomous agents.
A cross-thread observation from the Quality Engineering Safety
synthesis: items 1, 4, 7, and 10 share a common structural pattern —
the currently used ``healed / passed / approved'' signal is a single,
gameable proxy that is not independently verified — motivating a
general design principle for next-generation QE safety mechanisms:
\emph{require a second, independently-sourced signal before trusting
the first} (see Section~\ref{sec:qe-safety} and Section~\ref{sec:safety-framework}).


\section*{Funding}
This review did not receive any specific grant from funding agencies in the public, commercial, or not-for-profit sectors. The author conducted this study independently without external financial support.

\section*{Declaration of competing interest}
The author declares that they have no known competing financial interests or personal relationships that could have appeared to influence the work reported in this paper.

\section*{Declaration of generative AI and AI-assisted technologies in the manuscript preparation process}
During the preparation of this work, the author used AI-assisted tools for literature search assistance and text formatting. After using these tools, the author reviewed and edited the content as needed and takes full responsibility for the content of the publication.

\section*{Data availability}
The search log, screening records (inclusion/exclusion decisions), and per-study data-extraction forms that support the findings of this review are maintained in a structured project archive and are available from the author upon reasonable request. No new experimental data were generated for this review; all cited datasets and benchmarks (e.g., Defects4J, BugsInPy, SWE-bench) are publicly available from their original sources, referenced in the bibliography.

\printcredits

\section*{Acknowledgment}
The author thanks the anonymous reviewers for their constructive feedback.

\bibliographystyle{elsarticle-num}
\bibliography{references}

\end{document}
